\documentclass[11pt]{article}

\usepackage[margin=1in]{geometry}
\usepackage{amsmath,amssymb,amsthm,mathtools}
\usepackage{enumitem}
\usepackage{microtype}
\usepackage[hidelinks]{hyperref}
\usepackage{xurl}

\newtheorem{definition}{Definition}
\newtheorem{remark}{Remark}

\title{Primary-Path Separation Manifests for Anonymous-DHT Receipts\\
\large Guard Objects for Non-Collusion Assumptions (Series Synthesis)}
\author{}
\date{Unified archive note v0.05 (2026-02-27; archive v121)}

\begin{document}
\maketitle

\begin{abstract}
Many ``reader-privacy'' deployments rely on a primary path whose privacy story is not a pure numeric budget, but a \emph{separation assumption}: one role sees client-side metadata while a different role sees destination-side metadata, and the claim is interpreted only if these roles do not collude within the threat window.
This is a legitimate design axis (e.g., relay/gateway splits such as OHTTP\cite{rfc9458OHTTP}), but it is a common source of silent drift: the architecture name stays the same while the separation assumption quietly fails.

This note standardizes a small \emph{separation manifest}: a digest-bound guard object that makes the non-collusion assumption, its invalidation triggers, and its evidence hooks explicit and diffable across releases.
Other papers should cite this note instead of re-stating what ``separation class \texttt{sep.nr1}'' means.
\end{abstract}

\paragraph{Scope.}
This is an interface note.
It does \emph{not} claim that separation is sufficient for anonymity, and it does \emph{not} provide a universal accountant that turns separation assumptions into bits.
It only standardizes the guard object that makes such assumptions addressable, reviewable, and auditable under change control.

\paragraph{Imports.}
Receipt line items (including \texttt{obs\_model} and \texttt{knobs}) are owned by Synthesis~12.\cite{series:receiptitems}
Trace-interface vocabulary (exposure normal forms) is owned by Synthesis~3.\cite{series:traceifaces}
Threat-window declarations are owned by Synthesis~4.\cite{series:threatwindows}
Replay hooks and drift/equivalence semantics are owned by Synthesis~18--20.\cite{series:planstateequiv,series:auditorreplay}
User discovery/anchoring is owned by Synthesis~13 and Synthesis~9.\cite{series:userminiface,series:transparencyprimitives}
A concrete instantiation lives in the worked example (Synthesis~17).\cite{series:workedexample}

\paragraph{Exports.}
A minimal object shape for \emph{separation manifests} and a small editorial rule: \emph{separation-class claims must name a manifest digest, and manifest drift is interface drift unless explicitly coarsened.}
For deployments that want a small quantitative adjunct (separation failure as a receipt-grade $(0,\rho)$ checkpoint) without changing the receipt schema, see Synthesis~23.\cite{series:probsep}

\section{Why manifests (not prose)}
A ``separation'' claim fails in practice when a deployment accidentally re-introduces a linking handle between the separated roles:
shared keys, shared logs, co-resident operators, correlated identifiers, or a fallback that leaks destination information back to the client-side observer.
These failures are rarely visible from the headline protocol name.
A manifest makes them \emph{enumerable} and therefore diffable.

\section{Separation classes (names without semantics are not contracts)}
In this archive, a separation class (e.g., \texttt{sep.nr1}) is a human mnemonic for a particular non-collusion assumption.
The mnemonic is not a contract; the contract is a digest-bound manifest.

\begin{definition}[Separation manifest (minimal)]\label{def:sepmanifest}
A \emph{separation manifest} is a small, hash-addressable object containing:
\begin{enumerate}[leftmargin=*]
\item \textbf{Identity:} \texttt{sep\_manifest\_id} (digest id), plus a human label \texttt{sep\_label} (e.g., \texttt{sep.nr1}).
\item \textbf{Context binding:} the \texttt{tw\_id} (threat window) and the \texttt{exposure\_nf\_id} of the surface the manifest guards.\cite{series:threatwindows,series:traceifaces}
\item \textbf{Roles and partitions:} a list of roles (e.g., \texttt{client}, \texttt{relay}, \texttt{router}) and, for each role, a short description of which observation tiers/projections that role is assumed to see.
\item \textbf{Assignment / rotation policy:} how role instances are selected and rotated (e.g., per-request random relay selection from a declared set; per-epoch key rotation; explicit disjoint operator sets). This is the minimal hook needed to make separation claims replayable and (when desired) to justify a published separation-failure bound.\cite{series:probsep}
\item \textbf{Non-collusion clause:} an explicit set of forbidden collusion edges (e.g., \texttt{relay} $\not\leftrightarrow$ \texttt{router}) for the duration of the threat window.
\item \textbf{Linking handles:} the identifiers (keys, stable connection ids, log correlation ids) that would create linkability if shared across roles.
\item \textbf{Invalidation triggers:} a finite list of events/changes that cause the separation interpretation to become invalid (e.g., shared operator, shared key material, enabled debug logging, fallback policy change).
\item \textbf{Evidence hooks:} a replay-oriented pointer saying how an auditor would check the above (attestation type, log pointer, or a declared ``policy-only'' status).
\end{enumerate}
\end{definition}

\begin{remark}[Why this is receipt-grade]
A manifest is intentionally small.
It is not an evaluation trace; it is the \emph{guard surface} that explains what a ``0-bit'' primary claim is actually asserting.
It is therefore appropriate to treat it like a knob: commit to its digest and treat digest drift as claim drift.
\end{remark}

\section{How to carry the manifest without changing the schema}
Synthesis~12 deliberately keeps the public receipt schema minimal.
Separation manifests therefore live \emph{adjacent} to the receipt and are referenced by digest.
Two standard placements are:
\begin{itemize}[leftmargin=*]
\item \textbf{As a knob:} include \texttt{primary\_surface\_manifest\_digest} inside the line item's \texttt{knobs}.
\item \textbf{As an observation qualifier:} include \texttt{sep\_manifest\_id} inside the line item's \texttt{obs\_model}.
\end{itemize}
Either choice is acceptable as long as (i) the manifest digest is committed, and (ii) the compare profile surfaces it for drift checks.\cite{series:receiptitems,series:workedexample}

\section{Drift rules (editorial)}
A separation manifest is part of the interface claim.
Accordingly:
\begin{itemize}[leftmargin=*]
\item If \texttt{sep\_manifest\_id} changes, treat it as interface drift and issue a new receipt lineage (Certified~C + Release~A), unless the change is an explicit \emph{coarsening} that preserves the guarded surface.\cite{series:planstateequiv}
\item If the manifest invalidation triggers fire at runtime, the primary-path line item must be treated as out-of-scope (and any composite claim must fall back to the remaining budgeted surfaces).
\item If a deployment wants quantitative primary-path budgets, use Regime~II conditional certificates (Synthesis~21) and cite the spectral-to-budget translation (Anonymity~D) for delegation-based designs.\cite{series:condcert,series:spectral2budget}
\end{itemize}

\section{Research warranted}
This note standardizes the \emph{guard object}; it does not solve the ``how many bits does collusion probability buy me?'' problem.
Two promising directions are:
(i) probabilistic separation (explicitly budgeting collusion risk within a window) and
(ii) coupling guard manifests to measurable network-level independence evidence.
Both would benefit from being packaged as replayable, digest-bound evidence objects compatible with Synthesis~21.

Synthesis~23 records the minimal first step: how a published upper bound on separation failure can be interpreted as a guarded $(0,\rho)$ checkpoint in the series' existing endpoint semantics.\cite{series:probsep}

\begin{thebibliography}{99}\small

\bibitem{series:receiptitems}
Anonymous.
\newblock Receipt Line Items for Anonymity Claims: A Minimal Tuple Schema for Published Budgets.
\newblock Synthesis~12, The-One-True-Archive (2026).

\bibitem{series:traceifaces}
Anonymous.
\newblock Trace Interfaces for Anonymous DHTs.
\newblock Synthesis~3, The-One-True-Archive (2026).

\bibitem{series:threatwindows}
Anonymous.
\newblock Threat Models and Repetition Windows for Anonymous DHTs.
\newblock Synthesis~4, The-One-True-Archive (2026).

\bibitem{series:planstateequiv}
Anonymous.
\newblock Digest-Bound Replay Plans and State Declarations: Equivalence, Minimality, and Change Control.
\newblock Synthesis~18, The-One-True-Archive (2026).

\bibitem{series:auditorreplay}
Anonymous.
\newblock Auditor Replay for Anonymous-DHT Receipt Line Items: A Mechanical Checklist for Digest-Bound Claims.
\newblock Synthesis~20, The-One-True-Archive (2026).

\bibitem{series:userminiface}
Anonymous.
\newblock User-Verifiable Receipt Interfaces for Anonymous-DHT Claims.
\newblock Synthesis~13, The-One-True-Archive (2026).

\bibitem{series:transparencyprimitives}
Anonymous.
\newblock Transparency Primitives for Privacy Receipts: Beacons, VRFs, and Append-Only Logs.
\newblock Synthesis~9, The-One-True-Archive (2026).

\bibitem{series:workedexample}
Anonymous.
\newblock Worked Example: Instantiating a Tiered Reader-Privacy Receipt.
\newblock Synthesis~17, The-One-True-Archive (2026).

\bibitem{series:probsep}
Anonymous.
\newblock Probabilistic Separation Budgets: Separation Failure as a Receipt-Grade $(\varepsilon,\delta)$ Adjunct.
\newblock Synthesis~23, The-One-True-Archive (2026).

\bibitem{series:condcert}
Anonymous.
\newblock Conditional Per-Step Budget Certificates: A Receipt-Grade Interface for Adaptive Horizon Claims.
\newblock Synthesis~21, The-One-True-Archive (2026).

\bibitem{series:spectral2budget}
Anonymous.
\newblock Spectral Delegation for Anonymous DHT Lookups: From mixing certificates to receipt-grade per-step budgets.
\newblock Anonymity series Paper~D, The-One-True-Archive (2026).

\bibitem{rfc9458OHTTP}
M.~Thomson and C.~A.~Wood.
\newblock \emph{Oblivious HTTP}.
\newblock RFC 9458, IETF, January 2024.

\end{thebibliography}

\end{document}
